\documentclass[sfdefaults=false]{scrbook}

%\documentclass{scrreprt} %necessary to use to get two-sided title page

\KOMAoptions{fontsize=11pt, paper=a4, open=any} % Schriftgröße und Papierformat setzen.
\KOMAoptions{DIV=16} % Parameter mit dem man den Seitenrand ändern kann.

\KOMAoptions{toc=bibliography}

% \includeonly{%
%   thesis_intro,
%   thesis_appendix,
%   thesis_acknowledge
% }





% title page setup
\newcommand{\thesistitle}{Readout electronics for the BASiC experiment}
\newcommand*{\thesisauthor}{add author name}
\newcommand*{\thesistown}{Bonn}
%\renewcommand*{\InstituteName}{Helmholtz-Institut für Strahlen- und Kernphysik}

\newcommand*{\thesisrefereeonetext}{1. Gutachter}
\newcommand*{\thesisrefereeone}{Prof. Dr. GUTACHER EINS}
\newcommand*{\thesisrefereetwotext}{2. Gutachter}
\newcommand*{\thesisrefereetwo}{Dr. GUTACHTER ZWEI}

\newcommand*{\thesissubmit}{November 2026}
\newcommand*{\thesismonth}{November}
\newcommand*{\thesisyear}{2026}
\newcommand*{\thesisdedication}{hier könnte ihre werbung stehen}


\usepackage[onehalfspacing]{setspace} % sets 1.5 spacing

\usepackage{mathtools}
\usepackage{amssymb}
\usepackage{hyperref}
\usepackage[all]{hypcap}
\usepackage{subcaption}
\usepackage{float}
\usepackage[british]{babel}
\usepackage{siunitx}
\usepackage{physics}
\usepackage[version=4]{mhchem}
\usepackage{pdfpages}
\usepackage{longtable}
\usepackage{csquotes}

\captionsetup[subfigure]{justification=raggedright,singlelinecheck=off}

\sisetup{
	separate-uncertainty = true,
	multi-part-units=brackets,
	number-unit-product = \text{ },
	mode=math,
	exponent-product = \cdot,
	output-product = \cdot,
	output-decimal-marker = {.}
}

\usepackage{ circuitikz }

\usepackage[backend=biber, citestyle=numeric, bibstyle=numeric]{biblatex} % do not change this, the whole thing will break
% \usepackage{bibtex}
% \usepackage[backend=bibtex]{biblatex}

% \usepackage[headsepline=true,markcase=upper]{scrlayer-scrpage}
% \automark[chapter]{chapter}
% \ohead*{521: $\gamma$-Spektroskopie} % outer head on even & odd pages
% \setkomafont{pagehead}{\tiny\bfseries}


\addtokomafont{chapterprefix}{\raggedleft}
\addtokomafont{chapter}{\Large}
\addtokomafont{section}{\large}
\addtokomafont{subsection}{\large}
\addtokomafont{subsubsection}{\large}


\makeatletter
\renewcommand\chapter{\thispagestyle{plain}%
  \global\@topnum\z@
  \@afterindentfalse
  \secdef\@chapter\@schapter}
\makeatother

\DeclareSIUnit\angstrom{\protect \text  {Å}}
\DeclareSIUnit\frame{\protect \text{frame}}
\DeclareSIUnit\ms{\milli\s}
\DeclareSIUnit\mW{\milli\watt}
\DeclareSIUnit\nm{\nano\meter}
\DeclareSIUnit\M{\text{M}}
\DeclareSIUnit\rpm{\text{rpm}}


\usepackage{graphicx}


\DeclareSIUnit\bin{\text{Bin}}
\DeclareSIUnit\count{\text{Anzahl}}
\DeclareSIUnit\unity{}

%\titlehead{\centering\includegraphics[width=0.6\textwidth]{fig name here}} %for a title picture
\graphicspath{{figs/}}


\setlength\parindent{0pt}

% to check if labels are referenced
%\usepackage{refcheck}


% Glossary package
% \usepackage[acronym,toc,nosuper]{glossaries}

% \hypersetup{colorlinks=false,
%   linkbordercolor=blue,citebordercolor=magenta,urlbordercolor=darkgreen
% }

% Make a glossary and a list of acronyms.
% \makeglossaries
% \input{thesis_glossary}



% common macros
\def\diff#1#2{\frac{\text{d}#1}{\text{d}#2}}
\def\diffp#1#2{\frac{\partial#1}{\partial#2}}
\def\imply{\Longrightarrow}
\def\d#1{\text{d}#1\text{ }}
%\def\equiv{\Longleftrightarrow} % synonym for \iff

\def\OD{\text{OD}}


\def\todo#1{
	\textbf{\textcolor{blue}{[#1]}}
}



\def\fig#1{figure \ref{fig:#1}}
\def\tab#1{table \ref{tab:#1}}
\def\eq#1{equation \eqref{eq:#1}}

\def\jaft#1#2#3{
	\begin{table}[h]
	    \centering
	    \begin{tabular}{#1} \hline
		\input{#2}\end{tabular}
	    \caption{#3}
	    \label{tab:#2}
	\end{table}
}

\def\jaftl#1#2#3#4{
	\begin{table}[h]
	    \centering
	    \begin{tabular}{#1} \hline
		\input{#2}\end{tabular}
	    \caption{#3}
	    \label{tab:#4}
	\end{table}
}

\def\jafmt#1#2#3#4{
	\begin{table}[h]
	    \centering
	    \begin{tabular}{#1}
		#2
	    \end{tabular}
	    \caption{#3}
	    \label{tab:#4}
	\end{table}
}


\def\jafp#1#2{
	\begin{center}
		\includegraphics[width=0.8\textwidth]{#1}
		%\includegraphics[height=0.4\textheight]{#1}
		\captionof{figure}{#2}
		\label{fig:#1}
	\end{center}
}

\def\jafps#1#2#3{
	\begin{figure}[hbp]
		\centering
		\includegraphics[width=#3\textwidth]{#1}
		%\includegraphics[height=0.4\textheight]{#1}
		\captionof{figure}{#2}
		\label{fig:#1}
	\end{figure}
}

\def\jafpsw#1#2#3{
	\begin{figure*}[h]
		\centering
		\includegraphics[width=#3\textwidth]{#1}
		%\includegraphics[height=0.4\textheight]{#1}
		\captionof{figure}{#2}
		\label{fig:#1}
	\end{figure*}
}

\def\jafpp#1#2#3#4{
	\begin{figure}[htp]
		\centering
		\begin{subfigure}{0.45\textwidth}
		    \caption{}
			\label{fig:#1}

			\includegraphics[width = \textwidth]{#1}
			%\caption{}
			%\label{fig:#1}
		\end{subfigure}
		\begin{subfigure}{0.45\textwidth}
		    \caption{}
			\label{fig:#3}

			\includegraphics[width = \textwidth]{#3}
			%\caption{}
			%\label{fig:#3}
		\end{subfigure}
		\caption{a) #2 and b) #4}
	\end{figure}
}

\def\jafpsh#1#2#3{
	\begin{figure}[H]
		\centering
		\includegraphics[width=#3\textwidth]{#1}
		%\includegraphics[height=0.4\textheight]{#1}
		\captionof{figure}{#2}
		\label{fig:#1}
	\end{figure}
}

\def\ergo{\Rightarrow} %silly stuff
